%% kapa: exam=1 round=28 year=2017 period=1 form=A subject=civil_law subject_ko=민법 seq=9 number=9
%% source: Q-Net 감정평가사 기출문제 (cst00309) artlSeq=5199603 file=extracted/1008/5199603/2017년 제28회 감정평가사 1차시험 문제지/1교시 A형.pdf
%% conversion: kordoc 4.24.51 -> tools/gen1.py (auto); review=proofread-sonnet; flags=none
\begin{question}{9}
반사회질서 또는 불공정한 법률행위에 관한 설명으로 옳지 않은 것은? (다툼이 있으면 판례에 따름)
\begin{choices}
  \item 부첩관계의 종료를 해제조건으로 하는 증여계약은 그 조건뿐만 아니라 그 계약자체도 무효이다.
  \item 감정평가사를 통해 공무원에게 직무상 부정한 청탁을 하게 하고 그 대가로 상당한 금품을 교부하기로 한 약정은 무효이다.
  \item 불공정한 법률행위가 되기 위해서는 피해자의 궁박, 경솔, 무경험 중 어느 하나만 있으면 되고 그 모두가 있어야 할 필요는 없다.
  \item 계약이 불공정한 법률행위에 해당하여 무효라 하더라도 특별한 사정이 없는 한 그 계약에 관한 부제소합의는 유효하다.
  \item 법률행위의 내용이 반사회적인 것은 아니지만 반사회적 조건이 붙어 반사회적인 성질을 띠게 되면 그 법률행위는 무효이다.
\end{choices}
\end{question}
\answer{4}
